% Working notes on what we're doing and why.
% Raw material, not paper prose. Update as the project moves.

\section*{Process notes}

\subsection*{Where we are (2026-05-11)}

First extraction running on Mercury: Philadelphia DMA Spot
TV occurrences, 2014--2018, filtered to SSB-relevant Nielsen
PCC subcodes (F221/F222/F223 treated, F224 bottled water as
within-firm placebo). Single-DMA pass; no geographic control
yet. Pittsburgh + Baltimore + NYC + Boston + Cleveland are
queued as the next extraction.

\subsection*{What we want to learn from Philly}

Does at least one major SSB advertiser visibly cut Spot TV
buys in the Philadelphia DMA around the Jan 1, 2017 tax
(with anticipation possibly starting after the June 16,
2016 vote)? If the within-firm placebo (Coke vs.\ Dasani)
doesn't move and SSB spend does, that's the cleanest cut
this extraction can give. Geographic control comes in the
second extraction.

\subsection*{Reframing Keller, Guyt, Grewal (2024)}

Their three ``marketing conduct'' findings (promo frequency
$-2\%$, promo depth $-12\%$, feature frequency $-14\%$) net
out to a single object: effective pass-through is bigger
than the shelf-tag pass-through. Useful number, but it's a
decomposition of one price, not three separate margins.

For our purposes, the takeaway is: a manufacturer-side
advertising response, if it exists, is a genuinely separate
margin from anything in their paper. They look at retailer
in-store conduct on the demand side. We look at manufacturer
ad spend on the supply side. Different decision-maker,
different instrument, different mechanism.

\subsection*{Low pass-through is the precondition, not a
  side angle}

A manufacturer cuts ads after a tax only if its per-unit
margin actually falls. Anything else would be a sort-of second-order thing (you advertise less because people are more elastic). With full pass-through, consumers
eat the tax, margin is unchanged, no reason to touch the
ads FOC. With low pass-through, the firm absorbs the tax,
per-unit margin gets crushed, and the incentive to pull
back advertising is large. So pass-through heterogeneity
isn't a side angle --- it's the regime where any ad
response shows up at all, and the response should be
biggest exactly where pass-through is lowest.

The same logic applies in the sb\_incidence sportsbook
setting: the ZMC (zero marginal cost) regime is where
pass-through is low, the firm bears the tax, and the
advertising response should be biggest.

Practical implication for our cross-section:
\begin{itemize}
  \item If we look across SSB products and find some have
        $\sim$100\% pass-through and others much less, the
        ad response should concentrate in the low-pass-
        through group. Pooling across products dilutes the
        coefficient.
  \item For Philly specifically: published locality
        averages are 90--103\%. If those are accurate at
        the product level too, the margin hit is small and
        the predicted ad response is small. The interesting
        case is a sub-group with lower pass-through, if
        one exists.
  \item Boulder's 53--154\% range is huge --- some products
        eat way more of the tax than others. That's where
        product-level pass-through heterogeneity is most
        visible, even if Boulder is too small for Spot TV
        power.
\end{itemize}

\subsection*{What pass-through data exists}

\citet{keller_guyt_grewal_soda_marketing_conduct} report
locality-level rates (Philly 90--103\%, Boulder 53--154\%,
Oakland 92\%, SF 100\%, Seattle 59--90\%). These are
locality averages, not product-level. For the margin-
channel test we'd want \emph{within-locality, across-
product} pass-through heterogeneity. Possible sources:
\begin{itemize}
  \item Roberto et al.\ 2019 / Seiler, Tuchman, Yao 2021
        (Philly) -- check whether they report by
        brand/package size.
  \item Cawley et al.\ papers -- same check.
  \item We can also estimate it ourselves from Nielsen RMS
        scanner data at the UPC level for the same five
        localities.
\end{itemize}

\subsection*{Open questions}

\begin{itemize}
  \item Do any of the five localities have enough Spot TV
        spend to detect a DMA-month ad response at
        reasonable precision? Boulder almost certainly not.
        Philly probably yes. SF / Oakland / Seattle unclear.
  \item For DMAs that straddle the tax boundary (Oakland is
        in the SF DMA; the SSB tax applies to Oakland only),
        the DMA-level Spot TV measure averages over treated
        and untreated retail. That's a wash for our
        identification.
  \item Anticipation: most of these events have $\geq$ 6
        months between vote and implementation. Need to
        decide whether to estimate the effect starting from
        the vote date or the implementation date, and
        whether to report both.
\end{itemize}
